Tables~\ref{tab:preferred-clustering-comparison} and
\ref{tab:preferred-wild-bootstrap} extend the inference comparison while
holding coefficients and samples fixed. The first reports conventional tests
with sector-department intersections, sector, department, and two-way sector
and department clustering. The preferred and elasticity models contain
5,268,375 observations, 559 intersections, 17 sectors, and 33 departments.
The joint endpoint model contains 489,796 observations and 540 intersections,
with the same numbers of sectors and departments. The second table reports
wild cluster bootstrap tests.

Sector-department clustering allows dependence across workers and years
within each intersection. Sector clustering permits dependence across
departments within a sector, whereas department clustering permits dependence
across sectors within a department. Two-way clustering allows both, unlike
clustering on their intersections \citep{CameronMiller2015}.

Tests of individual coefficients and contrasts are two-sided. Their
\(p\)-values and 95 percent confidence intervals use a Student t distribution
with \(G-1\) degrees of freedom, where \(G\) is the number of clusters.
Two-way tests use the smaller cluster count, giving 16 degrees of freedom.
The variance of a contrast includes the covariances between its coefficients.
For joint hypotheses, we use Wald \(F\) tests. Their numerator degrees of
freedom equal the number of restrictions, and their denominator degrees of
freedom equal \(G-1\).

Because sector clustering leaves only 17 clusters, we also report
null-imposed wild cluster bootstrap tests following \citet{RoodmanEtAl2019}.
We use the WCR11 procedure with Rademacher weights and 9,999 draws for sector,
department, and two-way inference. In the two-way version, bootstrap weights
are drawn by sector and test statistics use the two-way covariance estimator.
These are two-sided, pointwise tests, without adjustment for multiple
comparisons. Coefficients, samples, and survey weights are held fixed across
inference methods. The bootstrap \(p\)-values supplement rather than replace
the baseline confidence intervals.

The formal 101+/solo premium is not significant with sector bootstrap
inference (\(p=0.184\)) or two-way bootstrap inference (\(p=0.188\)). By
contrast, every formal 101+ comparison with the categories from 2--3 through
51--100 remains significant at 5 percent. All eight informal size-to-solo
comparisons and all nine formal-informal wage gaps also remain significant
under each bootstrap convention. The tables display selected contrasts,
and the replication files report all 45 contrasts. The two-way bootstrap
\(p\)-value for the difference between the formal and informal 101+/solo
log-wage gaps is 0.092, compared with 0.0015 for the difference between their
elasticities. The 101+/solo contrast and the elasticity therefore provide
different evidence about how the size-pay relationship differs by formality.

The informal 2008--2025 increase has bootstrap \(p\)-values of 0.111 with
sector clustering, 0.266 with department clustering, and 0.273 with two-way
clustering. It is not significant under any of these bootstrap conventions.
The formal change is also insignificant. For that last comparison only,
two-way studentization produces 209 invalid bootstrap statistics, which the
software discards. Its reported \(p\)-value uses the remaining 9,790 draws.
The sector and department bootstrap tests of the formal change retain all
9,999 draws and give \(p=0.886\) and \(p=0.928\).

These checks address alternative assumptions about residual dependence,
not the exact GEIH sampling variance. GEIH uses stratified, multistage
geographical sampling, with successive selections ending in segments of
approximately ten dwellings \citep{DANE2008GEIHDesign,DANE2025GEIHDesign}.
The model file retains expansion weights but not the sampling-stratum,
primary-unit, or segment identifiers needed to reconstruct that design.
Sector-department intersections can separate workers who belong to the
same household or geographical sampling unit but work in different sectors.
Department and two-way clustering allow dependence across such sectors,
but neither reproduces the survey's stratification and selection stages.

\begin{table}[htbp]
  \centering
  \caption{Clustering sensitivity of preferred-model contrasts}
  \label{tab:preferred-clustering-comparison}
  \small
  \setlength{\tabcolsep}{5pt}
  \begin{tabular}{@{}p{0.40\textwidth}rrrrr@{}}
    \toprule
    Contrast & Estimate & $p$ (SD) & $p$ (S) & $p$ (D) & $p$ (S+D) \\
    \midrule
    \multicolumn{6}{l}{\textit{A. Within-formality log-wage gaps}} \\
    Formal 101+/solo & 0.153 & $<0.001$ & 0.1130 & $<0.001$ & 0.0900 \\
    Informal 101+/solo & 0.316 & $<0.001$ & $<0.001$ & $<0.001$ & $<0.001$ \\
    Formal minus informal, 101+/solo & -0.163 & $<0.001$ & 0.0603 & $<0.001$ & 0.0495 \\
    Formal 101+/2--3 & 0.225 & $<0.001$ & $<0.001$ & $<0.001$ & $<0.001$ \\
    Formal 101+/4--5 & 0.175 & $<0.001$ & $<0.001$ & $<0.001$ & $<0.001$ \\
    Formal 101+/6--10 & 0.140 & $<0.001$ & $<0.001$ & $<0.001$ & $<0.001$ \\
    Formal 101+/11--19 & 0.113 & $<0.001$ & 0.0015 & $<0.001$ & $<0.001$ \\
    Formal 101+/20--30 & 0.114 & $<0.001$ & $<0.001$ & $<0.001$ & $<0.001$ \\
    Formal 101+/31--50 & 0.093 & $<0.001$ & 0.0019 & $<0.001$ & $<0.001$ \\
    Formal 101+/51--100 & 0.067 & $<0.001$ & 0.0061 & $<0.001$ & 0.0024 \\
    Informal 101+/11--19 & -0.038 & 0.0241 & 0.2106 & 0.0027 & 0.1675 \\
    \midrule
    \multicolumn{6}{l}{\textit{B. Formal-minus-informal log-wage gaps}} \\
    Solo & 0.466 & $<0.001$ & $<0.001$ & $<0.001$ & $<0.001$ \\
    11--19 & 0.152 & $<0.001$ & $<0.001$ & $<0.001$ & $<0.001$ \\
    101+ & 0.303 & $<0.001$ & $<0.001$ & $<0.001$ & $<0.001$ \\
    \midrule
    \multicolumn{6}{l}{\textit{C. Change in the 101+/solo log-wage gap, 2008--2025}} \\
    Formal & 0.008 & 0.8609 & 0.8693 & 0.9074 & 0.9103 \\
    Informal & 0.152 & 0.0113 & 0.0896 & 0.1244 & 0.1987 \\
    \midrule
    \multicolumn{6}{l}{\textit{D. Approximate firm-size elasticities}} \\
    Formal & 0.040 & $<0.001$ & 0.0034 & $<0.001$ & 0.0022 \\
    Informal & 0.097 & $<0.001$ & $<0.001$ & $<0.001$ & $<0.001$ \\
    Formal-minus-informal & -0.057 & $<0.001$ & $<0.001$ & $<0.001$ & $<0.001$ \\
    \bottomrule
  \end{tabular}
  \vspace{0.3em}
  \begin{minipage}{0.98\textwidth}
    \footnotesize
    Notes: Panels A--B use Equation~(\ref{eq:main}), C uses the joint endpoint model in Equation~(\ref{eq:preferred-dynamics}), and D uses Equation~(\ref{eq:preferred-elasticity}). Estimates in A--C are in log points, not percentages or percentage-point changes. SD denotes sector-department intersections, S sector, D department, and S+D two-way sector and department clustering. Two-sided $t$ tests use 558 SD degrees of freedom (539 for C), 16 for S and S+D, and 32 for D.
  \end{minipage}
\end{table}

\begin{table}[htbp]
  \centering
  \caption{Wild cluster bootstrap tests for preferred-model contrasts}
  \label{tab:preferred-wild-bootstrap}
  \small
  \setlength{\tabcolsep}{5pt}
  \begin{tabular}{@{}p{0.40\textwidth}rrrr@{}}
    \toprule
    Contrast & Estimate & $p$ (S) & $p$ (D) & $p$ (S+D) \\
    \midrule
    \multicolumn{5}{l}{\textit{A. Within-formality log-wage gaps}} \\
    Formal 101+/solo & 0.153 & 0.1835 & $<0.001$ & 0.1883 \\
    Informal 101+/solo & 0.316 & $<0.001$ & $<0.001$ & $<0.001$ \\
    Formal minus informal, 101+/solo & -0.163 & 0.0859 & $<0.001$ & 0.0917 \\
    Formal 101+/2--3 & 0.225 & $<0.001$ & $<0.001$ & $<0.001$ \\
    Formal 101+/4--5 & 0.175 & 0.0015 & $<0.001$ & 0.0020 \\
    Formal 101+/6--10 & 0.140 & 0.0033 & $<0.001$ & 0.0037 \\
    Formal 101+/11--19 & 0.113 & 0.0035 & $<0.001$ & 0.0032 \\
    Formal 101+/20--30 & 0.114 & $<0.001$ & $<0.001$ & $<0.001$ \\
    Formal 101+/31--50 & 0.093 & 0.0059 & $<0.001$ & 0.0039 \\
    Formal 101+/51--100 & 0.067 & 0.0237 & $<0.001$ & 0.0169 \\
    Informal 101+/11--19 & -0.038 & 0.2165 & 0.0324 & 0.2095 \\
    \midrule
    \multicolumn{5}{l}{\textit{B. Formal-minus-informal log-wage gaps}} \\
    Solo & 0.466 & $<0.001$ & $<0.001$ & $<0.001$ \\
    11--19 & 0.152 & $<0.001$ & $<0.001$ & $<0.001$ \\
    101+ & 0.303 & $<0.001$ & $<0.001$ & $<0.001$ \\
    \midrule
    \multicolumn{5}{l}{\textit{C. Change in the 101+/solo log-wage gap, 2008--2025}} \\
    Formal & 0.008 & 0.8862 & 0.9284 & 0.9330$^{\dagger}$ \\
    Informal & 0.152 & 0.1113 & 0.2663 & 0.2728 \\
    \midrule
    \multicolumn{5}{l}{\textit{D. Approximate firm-size elasticities}} \\
    Formal & 0.040 & 0.0155 & $<0.001$ & 0.0173 \\
    Informal & 0.097 & $<0.001$ & $<0.001$ & $<0.001$ \\
    Formal-minus-informal & -0.057 & $<0.001$ & $<0.001$ & 0.0015 \\
    \bottomrule
  \end{tabular}
  \vspace{0.3em}
  \begin{minipage}{0.98\textwidth}
    \footnotesize
    Notes: Contrasts and specifications match Table~\ref{tab:preferred-clustering-comparison}. S and D denote sector and department clustering. S+D uses two-way covariance with bootstrap weights drawn by sector. Tests are two-sided WCR11 with the null imposed, Rademacher weights, and 9,999 draws. $\dagger$ The package discarded 209 invalid bootstrap $t$-statistics for this test, leaving 9,790 draws. All other displayed tests retained 9,999 valid draws. Tests are pointwise, without multiple-testing adjustment.
  \end{minipage}
\end{table}

\FloatBarrier
